
%----------------------------------------------------------------------------
%	Schedule
%----------------------------------------------------------------------------------------
\section{Schedule}
\subsection{Project Timeline}
The project timeline shows all major events that will be occurring as part of the STARS project. A flow chart is used in the figure below for the project timeline as it is a visual where one can see a clear sequence of events. The figure below shows the project timeline for NASC with some time estimations. 
\begin{figure}[H]
    \includegraphics[width=\columnwidth]{NASC_Critical_Path_PDR_Overview.png}\\
    \caption{NASC Complete Schedule}
    \label{fig:NASC Complete Schedule} 
\end{figure}  
The schedule is broken up into three distinct phases: Project Conceptual Phase, Design and Analysis Phase, and Ground Demonstration Phase. The Project Conceptual Phase is the early project phase where the NASC focused on setting the ground work for the project and combined the two originally separate teams, NASB and Wolf\^3. The next phase is the Design and Analysis Phases where the NASC's focus is towards developing the design and conducting analysis to validate the design for STARS. The Ground Demonstration Phase is where NASC will be focused on developing and testing different elements of the CubeSat design to validate models made during the design cycle.

The different colors represent the current status of the different tasks. The green tasks are tasks that have been completed at the writing of this document, yellow tasks represent currently in progress tasks, and red tasks represent future tasks. The blue boxes within the task boxes represent the approximate timeline for completion of the tasks. The critical path analysis conducted for PDR analysis estimate is found in section 8.2. The CDR analysis timeframe is an estimate based on current assumptions for PDR and is subject to change based on the timeframe for PDR analysis to be conducted. 
\subsection{PDR Analysis Critical Path}
Critical path analysis was used to determine the timeframe for the PDR analysis over alternative methods due to its ability to relate tasks and show how different tasks work in parallel with each other. This offers the advantage over standard Gantt charts where one can determine the critical path which drives the timeline for the overall project \cite{CPA}. The critical path for PDR Analysis can be found in the figure below.
\begin{figure}[H]
    \centering
    \includegraphics[width=\columnwidth]{NASC_PDR_Analysis.png}\\
    \caption{NASC PDR Analysis Critical Path}
    \label{fig:NASC PDR Critical Path} 
\end{figure}  
The different timetables for the PDR analysis were determined by estimating the time that each analysis was going to take and determining the dependency for the analysis to take place. For example the GNC subsystem is dependent on the thermal management subsystem because the thermal subsystem will have pointing requirements and maneuver requirements that feed into the GNC subsystem. The critical path is represented by a thicker red arrow which represents the longest sequence of events required to complete the PDR analysis. This critical path dictates the overall schedule for the analysis while other subsystems take reduced time and can be conducted in parallel. Other colored arrows are to make it easier to distinguish the different connections between subsystems. 